\documentclass[11pt]{article}
\usepackage[top=2.5cm,bottom=2.5cm,left=2.5cm,right=2.5cm]{geometry}
\usepackage{mathptmx} % Times New Roman-equivalent font
\usepackage[T1]{fontenc}
\usepackage[english]{babel}
\setlength{\parindent}{0pt}
\setlength{\parskip}{0pt}
\pagestyle{empty}

\begin{document}

\begin{center}
{\bfseries\fontsize{11}{13}\selectfont STATISTICAL DETECTION AND ATTRIBUTION OF AI-GENERATED AND HUMAN-WRITTEN PHISHING EMAILS USING NAIVE BAYES CLASSIFIERS}

\vspace{0.5cm}

{\bfseries Mannan Bhardwaj}

School of Computer Science Engineering and Information Systems, Department of Cyber Security, Vellore Institute of Technology, India

\vspace{0.3cm}

{\bfseries Bhoumik Singh Rajput}

School of Computer Science Engineering and Information Systems, Department of Cyber Security, Vellore Institute of Technology, India

\end{center}

\vspace{0.5cm}

\noindent{\bfseries ABSTRACT}

\vspace{0.2cm}

\noindent Phishing emails increasingly incorporate content generated by large language models, which blurs the linguistic boundary between traditional human-written phishing and legitimate correspondence. This study proposes a two-stage framework for distinguishing Legitimate, Human-Written Phishing, and AI-Generated Phishing emails using statistically validated linguistic and stylometric features rather than black-box deep learning. Fourteen quantifiable features spanning lexical, syntactic, and communication-style categories -- including token predictability, sentence-length variability, syntactic complexity, pronoun and urgency-keyword density, and lexical diversity -- are extracted from each email. Each feature is subjected to a rigorous, non-parametric statistical pipeline: Kruskal-Wallis H-tests and Chi-Square tests screen for genuine cross-class differences, Benjamini-Hochberg correction controls the false discovery rate across simultaneous tests, and effect-size measures (Cohen's d, Cram\'er's V) confirm practical significance before a feature is retained. Only statistically and practically validated features proceed to classification. Two specialized Naive Bayes sub-models -- a rank-binned model for continuous features and a categorical model for discrete features -- independently estimate class probabilities, which are then combined through a Meta-Naive Bayes layer into a single posterior distribution. Classification proceeds in two steps: first separating legitimate from phishing emails, then, for emails flagged as phishing, attributing authorship as human-written or AI-generated using a dedicated attribution-specific feature subset validated via Mann-Whitney U and Chi-Square tests. By grounding every retained feature in formal hypothesis testing, the proposed approach offers an interpretable, computationally lightweight alternative to transformer-based detectors, while remaining explicitly traceable back to measurable linguistic evidence behind each classification decision.

\vspace{0.3cm}

\noindent{\bfseries Keywords:} Phishing Detection, AI-Generated Text, Statistical Hypothesis Testing, Naive Bayes Classification, Stylometric Features

\end{document}
